\documentclass[10pt,a4paper]{amsart}
\usepackage[margin=19mm]{geometry}
\usepackage{amsmath,amssymb,mathtools}
\usepackage[scr=rsfs,cal=euler]{mathalfa}
\usepackage{xfrac}
\usepackage[unicode,hidelinks,bookmarks=false]{hyperref}
\newcommand{\ie}{i.e.\ }
\newcommand{\cf}{cf.\ }
\newcommand{\resp}{respectively\ }
\newcommand{\Subsection}{\subsection}
\providecommand{\nobreakdash}{\nobreak\hbox{-}}
\setlength{\emergencystretch}{2em}
\multlinegap0pt
\usepackage{bm}
\usepackage{bbm}
\usepackage{dsfont}
\usepackage{xspace}
\usepackage{cancel}
\usepackage{mathtools}
\usepackage{cleveref}
\usepackage[shortlabels]{enumitem}
\usepackage{amsthm}
\makeatletter
\@ifundefined{theorem}{\newtheorem{theorem}{Theorem}}{}
\@ifundefined{corollary}{\newtheorem{corollary}[theorem]{Corollary}}{}
\@ifundefined{definition}{\newtheorem{definition}[theorem]{Definition}}{}
\@ifundefined{lemma}{\newtheorem{lemma}[theorem]{Lemma}}{}
\@ifundefined{proposition}{\newtheorem{proposition}[theorem]{Proposition}}{}
\@ifundefined{remark}{\newtheorem{remark}[theorem]{Remark}}{}
\makeatother
\def\C{\mathbb{C}}
\def\DC{\textsf{DC}}
\def\GL{\textsf{GL}}
\def\Id{\textsf{Id}}
\def\N{\mathbb{N}}
\def\O{\mathcal{O}}
\def\Q{\mathbb{Q}}
\def\S{\textsf{QuickWP}}
\def\WP{\textsf{WP}}
\def\Z{\mathbb{Z}}
\def\mat{\textsf{M}}
\def\tH{\tilde{H}}
\def\tU{\tilde{\mathfrak{U}}}
\title[The average-case complexity of the Word Problem for groups o]{The average-case complexity of the Word Problem for groups of matrices over $\mathbb{Z}$ is linear}
\author{Fr\'ed\'erique Bassino, Cyril Nicaud, Pascal Weil}
\date{Source: arXiv:2506.00948v2 (CC BY 4.0)}
\begin{document}
\maketitle

\noindent\emph{Source and licence.} The average-case complexity of the Word Problem for groups of matrices over $\mathbb{Z}$ is linear. Version arXiv:2506.00948v2 (CC BY 4.0), distributed under the Creative Commons Attribution 4.0 International licence, as recorded in the arXiv licence metadata for this exact version and shown on the abstract page https://arxiv.org/abs/2506.00948v2. Full rights notice: licenses/arxiv-2506.00948-CC-BY-4.0.txt.\par\smallskip

\noindent\textit{Editorial scope.} Counted unit: the Theorem whose source label is \texttt{thm: main for integer rings} in the article (source0.tex lines 760--764, source label \texttt{thm: main for integer rings}), delivered together with the article's own complete original proof of it (source0.tex lines 766--794, one \texttt{proof} environment of 29 lines). No mathematics is written, repaired or re-derived: statement and proof are the article's own text line for line. Required context retained uncounted: prop: basic facts (lines 226--246); rk: worst case in other rings (lines 375--377); prop: upper triangular (lines 390--394); def: def of q (lines 408--413); lemma: ac of DCm (lines 463--469); thm: technical theorem (lines 489--493); lm: matrices of small order (lines 640--642); cor:convergence uniform lazy (lines 679--685). Every definition, display and label of those lines is retained unchanged. Omitted from this package and not redistributed: 1--225, 247--374, 378--389, 395--407, 414--462, 470--488, 494--639, 643--678, 686--759, 765--765, 795--945. Nothing inside a retained excerpt is omitted or altered: every retained body is delivered exactly as the article's lines are written, so no omission receipt of any kind accompanies this package. Citations: the retained excerpts carry no citation command at all, so no bibliography entry of the article is transcribed and no reference is redistributed. Reference adaptation: none was needed and none was made; every reference inside the delivered unit resolves inside it, so no unresolved reference or citation is produced. Printed slips kept literal: no printed word, symbol, hypothesis or number of the counted statement or of its proof was altered. Layout and class: the article's own class is \texttt{article} with packages including graphicx, a4wide, amsmath, amssymb, amsthm, xcolor, hyperref, paralist, cleveref, comment, algorithm2e; the wrapper is the lane's pinned \texttt{amsart} editorial wrapper at 10pt on A4, which restates the theorem-like environments the retained text needs and adds the article's own macro definitions that the retained text needs (\texttt{\textbackslash C}, \texttt{\textbackslash DC}, \texttt{\textbackslash GL}, \texttt{\textbackslash Id}, \texttt{\textbackslash N}, \texttt{\textbackslash O}, \texttt{\textbackslash Q}, \texttt{\textbackslash S}, \texttt{\textbackslash WP}, \texttt{\textbackslash Z}, \texttt{\textbackslash mat}, \texttt{\textbackslash tH}, \texttt{\textbackslash tU}), with extra wrapper packages (bm, bbm, dsfont, xspace, cancel, mathtools, cleveref, enumitem). The delivered numbering is the unit's own. Assets: no retained span contains a figure environment, a graphics-inclusion command, a PDF-page inclusion, a table or a TikZ picture; no image file is copied into this package, the package declares no asset dependency and no image file is redistributed with it. Licence: version arXiv:2506.00948v2 of this article is distributed under the Creative Commons Attribution 4.0 International licence, as recorded in the arXiv licence metadata for this exact version and shown on the abstract page https://arxiv.org/abs/2506.00948v2. A full rights notice is delivered at licenses/arxiv-2506.00948-CC-BY-4.0.txt. Characters: every retained excerpt is pure ASCII, so the delivered engine, XeLaTeX with its default Latin Modern text font, reports no missing character.
\begin{proposition}\label{prop: basic facts}
%
Let $n,n'$ be integers and let $L$ such that $\ell(n), \ell(n') \le L$.
%
\begin{enumerate}[(i)]
%
\item $\ell(nn') \le \ell(n) + \ell(n') - 1$. 
%
\item\label{fact: HvdH} The product $nn'$ is computed in time $\O(L\log L)$.

\item The product $nn'$ is also computed (by the primary school multiplication algorithm) in time $\O(\ell(n)\ell(n'))$, which is interesting if $\ell(n)$ is small with respect to $\ell(n')$.

\item The sum of $d$ integers of length at most $L$, has length at most $1 + \log d + L$ and is computed in time $\O(d\log d + dL)$.

\item\label{fact: product by general matrix} If $A, A'$ are $d \times d$ matrices with entries of length at most $L$, the entries of the product $AA'$ have length at most $1 + \log d + 2L$, and each is computed in time $\O(d\log d + dL\log L)$.

\item\label{fact: product by fixed matrix} If the entries of $A$ (resp. $A'$) are of length at most $L$ (resp. $L'$), and if $L'$ is much smaller than $L$, the entries of the product $AA'$ have length at most $1 + \log d + L + L'$, and each is computed in time $\O(d\log d + dLL')$.
%
\end{enumerate}
%
\end{proposition}

\begin{remark}\label{rk: worst case in other rings}
    In fact, the same algorithm runs on matrices over any computable ring. Over $\Q$, it also yields an $\O(n\log^2n)$ worst-case complexity since the addition and multiplication of rationals takes asymptotically the same time as the addition and multiplication of integers.
\end{remark}

\begin{proposition}\label{prop: upper triangular}
    If $\Sigma$ consists only of upper (resp. lower) triangular matrices, then the worst-case complexity of $\WP_\Sigma$ is $\O(n)$.

    If $G$ is a finitely generated nilpotent group, then the Word Problem in $G$ can be solved in linear worst-case complexity.
\end{proposition}

\begin{definition}\label{def: def of q}
    Let $q\colon \N\to\N$ be the function given by $q(0) = q(1) = 1$ and, for all $n\ge 2$,
\[
q(n) = \prod_{\substack{p \enspace \leq \enspace \log^5 n\\p\text{ prime}}} p
\]
\end{definition}

\begin{lemma}\label{lemma: ac of DCm}
%
Let $P_n$ be the probability that a word $w$ of length $n$ satisfies $\mat(w)_{q(n)} = \Id$. The aver\-age-case complexity of $\S$ (when inputs are uniform random words of length $n$) is $\O(n + P_n\, n\, \log^2n)$.

Moreover if $H = \langle\Sigma\rangle$ is finite, then the average-case complexity of $\S$ is $\O(n)$.
%
\end{lemma}

\begin{theorem}\label{thm: technical theorem}
    %
    If $H = \langle\Sigma\rangle$ is infinite, then $\mat(w)_{q(n)} = \Id$ with probability $\O(\log^{-2} n)$ (when $w$ is chosen uniformly at random among length $n$ words).
    %
\end{theorem}

\begin{lemma}\label{lm: matrices of small order}
Let $A\in\GL_d(\mathbb{Z})$ be a matrix of infinite order.  The number of primes $p$ such that $A_p$ has order at most $\mathfrak{o}$ in $\GL_d(\Z/p\Z)$ is $\O(\mathfrak{o}^2)$. 
\end{lemma}

\begin{corollary}\label{cor:convergence uniform lazy}
%
Let $q(n)$ be the map defined in \cref{def: def of q}, and let $(h_n)_n$ be a sequence such that $h_n \in \tH_{q(n)}$ for each $n$. 
%
If $H$ is infinite, then the probability that a length $n$ trajectory in $\tU_{q(n)}$ ends in $h_n$ is $\O(\log^{-2}n)$.
%
\end{corollary}

\begin{theorem}\label{thm: main for integer rings}
    Let $R$ be a subring of $\C$ such that the additive group $(R,+)$ is finitely generated. The Word Problem $\WP_\Sigma$ in $\GL_d(R)$ has worst-case complexity $\O(n\log^2n)$ if $\langle\Sigma\rangle$ is infinite, and $\O(n)$ if $\langle\Sigma\rangle$ is finite, or if $\Sigma$ consists of upper triangular matrices.
    
    Moreover Algorithm $\S$ solves $\WP_\Sigma$ in $\GL_d(R)$ with linear average-case complexity.
\end{theorem}

\begin{proof}
%
The proof is essentially the same as for matrices over $\Z$. The following is a discussion of the necessary adjustments of the reasoning.

If $x = \sum_{i=1}^\delta x_i\alpha_i$ is an element of $R$ (where each $x_i$ is in $\Z$), we let $\|x\|_\infty = \max_i|x_i|$. If $A\in \GL_d(R)$, we let $\|A\|_\infty = \max\{\|a_{i,j}\|_\infty \mid 1\le i,j\le d\}$. Finally, we let $\beta = \max\{\|\alpha_i\alpha_j\|_\infty \mid 1\le i,j \le \delta\}$. We also let $\ell(x) = \ell(\|x\|_\infty)$ and $\ell(A) = \ell(\|A\|_\infty)$.
%Recall that, for any integer $n$, $\ell(n) = \lceil\log(|n|+1)\rceil+1$ and $|n| \le 2^{\ell(n)}$, see \cref{sec: computing with integers}.
%
The following facts are verified using \cref{prop: basic facts}.
%
\begin{itemize}
    \item If $x_1,\dots, x_k \in R$ with each $\ell(x_i) \le L$, then $\sum_i x_i$ is computed in time $\O(L)$ and $\|\sum_ix_i\|_\infty \le k \max(\|x_i\|_\infty)$.

    \item If $x, y \in R$ and $\ell(x), \ell(y) \le L$, then $xy$ is computed in time $\O(L\log L)$ and $\|xy\|_\infty \le \beta\, \|x\|_\infty \|y\|_\infty$.

    \item If $A, B \in \GL_d(R)$ and $\ell(A), \ell(B) \le L$, then $\|AB\|_\infty \le d\beta\, \|A\|_\infty \|B\|_\infty$, and it is computed in time $\O(L\log L)$.

    \item Suppose now that $y \in R$ is a constant and $x\in R$. Then $\|xy\|_\infty = \O(\|x\|_\infty)$, and $xy$ is computed in time $\O(\ell(x))$. If $B$ is a constant matrix and $A \in \GL_d(R)$, the size of the coefficients of $AB$ is $\O(\|A\|_\infty)$, and $AB$ is computed in time $\O(\ell(A))$.
%
\end{itemize}

The facts listed above, along with \cref{rk: worst case in other rings}, show that the worst case complexity of $\DC_\Sigma$ (on $\GL_d(R)$) is $\O(n\log^2n)$ if $H$ is infinite, and $\O(n)$ if $H$ is finite. The reasoning in \cref{prop: upper triangular} also applies in this case: if $\Sigma$ consists only of upper triangular matrices, then the worst case complexity of $\DC_\Sigma$ is $\O(n)$.

As for matrices over $\Z$, \cref{lemma: ac of DCm} applies and it reduces the proof to proving an analogue of \cref{thm: technical theorem}: verifying that, if $H = \langle\Sigma\rangle$ is infinite, then $\mat(w)_{q(n)} = \Id$ with probability $\O(\log^{-2} n)$ (when $w$ is chosen uniformly at random among length $n$ words).

A minor remark concerns the proof of \cref{lm: matrices of small order} when $\Z$ is replaced by $R$. We note that, if $A$ is a matrix of infinite order with coefficients in $R$, then $\|A^n\|_\infty \le (d\beta)^{n-1}\|A\|_\infty^n$ (instead of $d^{n-1}\|A\|_\infty^n$). This is no obstacle in the rest of the proof of that lemma.

Similarly, the proof of \cref{cor:convergence uniform lazy} carries as in the case of coefficients in $\Z$, since Schur's theorem applies to all finitely generated subgroups of $\GL_d(\C)$, such as $H = \langle\Sigma\rangle$.
%
\end{proof}

\end{document}
